\ssr{РЕФЕРАТ}

Расчетно-пояснительная записка по курсовой работе на тему <<Разработка программного решения для трехмерной визуализации и генерации городской среды>> содержит~\pageref{LastPage} с., \totalfigures\ рис., 1 табл., \totalenumivs\ ист, 4 прил. 

Работа посвящена разработке программного комплекса для процедурной генерации и интерактивной 3D-визуализации городской среды. Цель достигнута через реализацию алгоритмов построения связной транспортной сети с множеством перекрёстков, создания моделей зданий, а также освещения на основе карты теней. Обеспечены воспроизводимость сцен и интерактивное управление камерой.  

В основе решения --- применение Z-буфера, диффузной закраски и аффинных преобразований. Приложение написано на C++20 и Qt6.5 и разделено на модули генерации, рендеринга и UI. Исследования показали квадратичную зависимость времени отрисовки от количества пикселей кадра.  

Ключевые слова: Z-буфер, город, карта теней, C++20, Qt6.5.